\chapter{Conclusion}
\label{chap:conclusion}

\begin{sloppypar}
This qualification develops one scientific argument. Rebalancing and unconstrained augmentation do not necessarily resolve weak prediction when outcome classes overlap. Domain-guided filtering can help, but its effects depend on the anchor and classifier. This motivated a separability diagnosis and then SepAware, an explicit intervention on class-conditional geometry. The factorial experiment attributes the strong Vigitel TSTR gain to separability pressure: it explains 96.00\% of treatment-plus-error variation, whereas anchor and interaction contributions are negligible. The corresponding COVID-19 effect is not established; separability explains 5.13\% and residual error 81.53\%.

The current contribution is consequently a method, a component-level empirical finding, and a set of boundary conditions---not proof of universal superiority. SepAware may still select easy prototypes, reduce minority diversity, amplify shortcuts, or increase privacy risk. The remaining programme tests those explanations by separating imbalance from overlap, repeating comparisons across generators and seeds, measuring clinical, causal, and privacy consistency, and validating under external shift. The intended PhD contribution is a defensible account of when separability should be controlled, how it trades against fidelity and safety, and when it should not be used.
\end{sloppypar}
